%%%%%%%%%%%%%%%%%%%%%%%%%%%%%%%%%%%%%%%%%%%%%%%%%%%%%%%%%%%%%%%%%%%%%%%%%%%%%%%%%%%%%%%%%%
% Bachelorarbeit/Projektarbeit/Masterarbeit
% Lehrstuhl für Strömungsmechanik (LSTM)
% Friedrich-Alexander-Universität (FAU) Erlangen-Nürnberg
%%%%%%%%%%%%%%%%%%%%%%%%%%%%%%%%%%%%%%%%%%%%%%%%%%%%%%%%%%%%%%%%%%%%%%%%%%%%%%%%%%%%%%%%%%

\newif\ifnomenklatur
\nomenklaturfalse
% Layout des Dokuments
\documentclass[
% a4paper,             					% Benutzt DIN-A4-Papier-Format (=default)
 	11pt,                 				% Schriftgröße 11 Punkte
  	notitlepage,          				% Kein Titelblatt wird erstellt
  	twoside, 			   				% Es werden Vorder- und Rückseite der Blätter benutzt
 	% openright,           				% Kapitel beginnen auf rechter Seite
	% onecolumn,           				% Der Text ist einspaltig (=default)
 	% bibtotoc,            				% Das Literaturverzeichnis wird als nummeriertes Kapitel in das Inhaltsverzeichnis aufgenommen
 	bibliography=totoc,
 	%liststotoc,          				% Tabellen- und Abbildungsverzeichnisse werden als Kapitel in das Inhaltsverzeichnis aufgenommen
  	listof=totoc,
  	headsepline,         				% Die Kopfzeile wird durch eine horizontale Linie vom Text getrennt
  	footsepline,         				% Die Fußzeile wird durch eine horizontale Linie vom Text getrennt
%  	chapterprefix,       				%	 Über den Kapitelüberschriften steht immer das Wort "Kapitel" mit der entsprechenden Nummer
  	%pointlessnumbers     				% Die Abschnittsnummerierungen werden immer ohne den letzten Punkt dargestellt, also "1.2.3" statt "1.2.3."
  cleardoublepage=empty,				% leere Füllseiten sind leer
  numbers=noenddot,
  %draft
]{scrreprt}


\setkomafont{sectioning}{\rmfamily\bfseries} % Überschriften mit Serifen


% \ifenglish
\usepackage[ngerman,english]{babel}	    		% Überschriften werden von LaTeX in der korrekten Sprache erzeugt (english = Hauptsprache)
% \else
% \usepackage[english,ngerman]{babel}	    		% Überschriften werden von LaTeX in der korrekten Sprache erzeugt (ngerman = Hauptsprache)
% \fi

%\usepackage[ansinew]{inputenc}			% Umlaute
%\usepackage[utf8]{inputenc}
%\usepackage{subfig}
\usepackage[font=small,labelfont=bf]{subcaption}
\usepackage{placeins}
\usepackage{preview}
\newcommand{\kreis}[1]{\unitlength1ex\begin{picture}(2.5,2.5)\put(0.75,0.75){\circle{2.5}}\put(0.75,0.75){\makebox(0,0){#1}}\end{picture}}
\usepackage{diagbox}
\usepackage{colortbl}
\usepackage[T1]{fontenc}		        % Schriftart
\usepackage[utf8]{inputenc}	  		% Direkte Eingabe der Sonderzeichen über die Tastatur 
%\usepackage{lmodern}					% schönere Schrift
%\usepackage{ae}							% schöne Schrift
%\usepackage[final]{microtype}			% mikrotypografische Schriftanpassung
\usepackage{lipsum}
%\DisableLigatures{}

%\usepackage{fixltx2e}
\usepackage{graphicx}	
\usepackage{epstopdf}	
\usepackage{pgf}	       		% Ermöglicht die Einbindung von Bildern
\providecommand{\mathdefault}[1]{#1}	% benötigt für matplotlib pgf-Export
\usepackage{makeidx}		          	% Erzeugt Verzeichnisse
\usepackage{array}
\usepackage{float}		             	% Erzwingt Positionierung von Floatobjekten mit "`H"'
\usepackage{tabularx}		          	% Tabellen mit Spalten gleicher Breite
%\usepackage{tabu}						% sehr flexible Tabellen mit mehr Möglichkeiten
%\usepackage{dcolumn}		          	% Tabellen mit ausgerichteten Einträgen
\usepackage{amssymb,amsmath}	    	% Um Texte in mathematischen Gleichungen zu schreiben mit dem Befehl \text{}
%\usepackage[figuresright]{rotating}	% Damit können Bilder, Tabellen, usw. in landscape Format angezeigt werden
\usepackage{bibgerm}		           	% Ermöglicht das Einbinden deutschsprachiger Literaturverzeichnisse
%\usepackage{subfig}		            % Mehrere Bilder nebeneinander, unter gleicher Bildunterschrift und -nummer
\usepackage{booktabs}		          	% Hochwertige Tabellen
\usepackage{longtable}	        		% Tabellen über mehrere Seiten
\usepackage{url}			            % Weblinks
\usepackage{pdfpages}					% PDF einbinden
%\usepackage{eurosans}
\usepackage[bookmarks,pagebackref=false,pdftex,bookmarksopen=true,bookmarksnumbered]{hyperref}	% PDF-Bookmarks, Angabe im Literaturverzeichnis auf welcher Seite Einordnungsformel steht momentan auf false gesetzt
\usepackage{multirow}
%\usepackage[square]{natbib}         	% nützliches Paket, falls man die Gestaltung des Literaturverzeichnisse und der Referenzen selbst übernehmen möchte
\usepackage{wrapfig}
\usepackage{tocvsec2}
\usepackage{ifthen}
\usepackage{xspace}						%Leerzeichen nach Anführungszeichen
%\usepackage[decimalsymbol=comma, expproduct=cdot]{siunitx}				%Einheiten im Mathemodus
%\usepackage{microtype,textcomp}			% bei siunitx
\usepackage{bm}
\usepackage{pgfplots}					%Matlab Diagramme in Latex
\def\rb#1{\rotatebox{90}{$\xleftarrow{#1}$}}
\newcommand{\tikzmark}[1]{\tikz[overlay, remember picture] \coordinate (#1);}


%---------------------------------------
% Glossary
%---------------------------------------
\usepackage[acronym, toc, nonumberlist]{glossaries}
\makeglossaries



\graphicspath{{./plots/}}


\usetikzlibrary{plotmarks}
\usepackage{cancel}				%Terme durchstreichen
\usetikzlibrary{plotmarks}
\usepackage{pgfplots}
\pgfplotsset{compat=1.4, every axis legend/.style={
	y tick label style={/pgf/number format/1000 sep=},					
	x tick label style={/pgf/number format/1000 sep=},},
	legend style={font=\scriptsize}}
	\newlength\figureheight  \newlength\figurewidth 
\usepackage{nicefrac} 										% erstellt schräggestelle Brüche
\tikzset{every picture/.style={line width=1pt}}				% Linewidth global für Tikz
\tikzset{every picture/.style={mark size=2}}				% Markersize global für Tikz
	
\usepackage{array}

\hyphenation{non-di-men-sio-na-lize non-di-men-sio-na-li-za-tion}



\usepackage[round]{natbib} 

%%%%%%%%%%%%%%%%%%%%%%%%%%%%%%%%%%%%%%%%%%%%%%%%%%%%%%%%%%%%%%%
% folgende Pakete für tikz-Bilder
%%%%%%%%%%%%%%%%%%%%%%%%%%%%%%%%%%%%%%%%%%%%%%%%%%%%%%%%%%%%%%%
\usepackage{tikz}
\usepackage{tikz-3dplot}
\usetikzlibrary{shapes,arrows,patterns,decorations.markings}
%\usepackage{color}
\usepackage[margin=0cm,nohead]{geometry}
% needed for BB
\usetikzlibrary{calc}
\usetikzlibrary{patterns}
% alle Pfeile global ändern

\tikzset{>=latex}
% Strich-Punkt-Linie
\tikzset{dashdot/.style={dash pattern=on .8pt off 3pt on 10pt off 3pt}}	
%Kreuz		
\tikzset{cross/.style={cross out, draw=black, fill=none, minimum size=4pt, inner sep=0pt, outer sep=0pt}, cross/.default=2.5pt}
%Strichstärken einstellen															
\tikzstyle{thin} = [line width=0.5pt]
\tikzstyle{semithick} = [line width=0.7pt]
% Manuelles Pattern erstellen
\pgfdeclarepatternformonly[\LineSpace]{my north east lines}{\pgfqpoint{-1pt}{-1pt}}{\pgfqpoint{\LineSpace}{\LineSpace}}{\pgfqpoint{\LineSpace}{\LineSpace}}%
{
	\pgfsetlinewidth{1pt}
	\pgfpathmoveto{\pgfqpoint{0pt}{0pt}}
	\pgfpathlineto{\pgfqpoint{\LineSpace + 0.1pt}{\LineSpace + 0.1pt}}
	\pgfusepath{stroke}
}
\newdimen\LineSpace
\tikzset{
	line space/.code={\LineSpace=#1},
	line space=3pt
}


%%%%%%%%%%%%%%%%%%%%%%%%%%%%%%%%%%%%%%%%%%%%%%%%%%%%%%%%%%%%%%%%%%%%%%%
 % Nomenklatur erstellen
\usepackage[norefeq,intoc]{nomencl}
\setlength{\nomitemsep}{-0.5\parsep}
\makenomenclature
\usepackage{etoolbox}
\makeindex

\renewcommand{\nomgroup}[1]{\vspace{0.5cm}
	\ifthenelse{\equal{#1}{G}}{\item[\textbf{Greek letters}]}{%
		\ifthenelse{\equal{#1}{L}}{\item[\textbf{Latin letters}]}{%
			\ifthenelse{\equal{#1}{I}}{\item[\textbf{Indexes}]}{%
				\ifthenelse{\equal{#1}{A}}{\item[\textbf{Abbreviations}]}{%
					\ifthenelse{\equal{#1}{O}}{\item[\textbf{Operators}]}{}
				}
			}
		}
	}
	\hspace*{-\leftmargin}\vspace{0.25cm} 
}


%%%%%%%%%%%%%%%%%%%%%%%%%%%%%%%%%%%%%%%%%%%%%%%%%%%%%%%%%%%%%%%%%%

%%%%%%%%%%%%%%%%%%%%%%%%%%%%%%%%%%%%%%%%%%%%%%%%%%%%%%%%%%%%%%%%%%
% PDF-Erstellung
%%%%%%%%%%%%%%%%%%%%%%%%%%%%%%%%%%%%%%%%%%%%%%%%%%%%%%%%%%%%%%%%%%
\hypersetup{pdftitle = {Direct numerical simulation of turbulent flow in rough pipes},
pdfauthor = {Jonas Tiefenthäler}}

\hypersetup{
	pdfstartview=FitV,
	bookmarksopen=true,
	bookmarksnumbered=true,
	breaklinks=true,
   	colorlinks = true,
   	linkcolor = black,
   	anchorcolor = black,
   	citecolor = black,
   	filecolor = black,
   	pagecolor = black,
   	urlcolor = black}


\newcommand{\comment}[1]{}		          % Um einen ganzen Textblock auszukommentieren 
\newcommand{\Rey}{\mathrm{Re}}   		  % Reynoldszahl
\renewcommand{\baselinestretch}{1.15}	  % Zeilenabstand 115%
\setlength{\belowcaptionskip}{5pt}	  % Abstand von Bildunterschriften zum Bild?
\setlength{\abovecaptionskip}{8pt}
\pagestyle{headings}		              % Die Kapitelüberschrift wird in der Kopfzeile angezeigt
%\frenchspacing                  	      % Kein Mehrabstand nach Satzende
\setcounter{tocdepth}{5}
\setcounter{secnumdepth}{5}		          % Kapitelnummerierung nur bis zur Ebene "subsection" durchgeführt
\let\endgraph\endgraf			          % falls nicht die neueste Version installiert ist, würde endgraf eine Fehlermeldung produzieren

%%%%%%%%%%%%%%%%%%%%%%%%%%%%%%%%%%%%%%%%%%%%%%%%%%%%%%%%%%%%%%%%%%%%%%%%%%%%%%%%%
% Bestimmung des Textfeldes
%%%%%%%%%%%%%%%%%%%%%%%%%%%%%%%%%%%%%%%%%%%%%%%%%%%%%%%%%%%%%%%%%%%%%%%%%%%%%%%%%
% Die Vorlage ist so gedacht, dass der Text doppelseitig gedruckt wird!
\setlength{\topmargin}{-15mm}	   	% Abstand zwischen dem oberen Rand jeder logischen Seite und der Oberkante der Kopfzeile.
\setlength{\headheight}{11mm}	   	% Höhe Kopfzeile
\setlength{\headsep}{7mm}	       	% Abstand zwischen Unterkante Kopfzeile und Rumpf
\setlength{\textheight}{237mm}	 	% Definiert die GesamtHöhe des Textrumpfes für alle nachfolgenden Seiten.
\setlength{\footskip}{10mm}	     	% Abstand zwischen Rumpf der Seite und Unterkante Fußzeile
\setlength{\topskip}{0mm}	       	% Abstand fest, der am oberen Rand des Seitenrumpfes zusätzlich eingefügt wird.
\setlength{\oddsidemargin}{+4mm} 	% bei zweiseitiger Format. den linken Rand der Seiten mit ungerader Seitennummer.
\setlength{\evensidemargin}{-4mm} 	% bei zweiseitiger Format. linken Randes der Seiten mit gerader Seitennummer.
\setlength{\textwidth}{160mm}	   	% Definiert die GesamtBreite des Textrumpfes für alle nachfolgenden Seiten.
\setlength{\parindent}{0mm}	   		% keine Einrückung bei neuem Absatz
%%%%%%%%%%%%%%%%%%%%%%%%%%%%%%%%%%%%%%%%%%%%%%%%%%%%%%%%%%%%%%%%%%%%%%%%%%%%%%%%%

%%%%%%%%%%%%%%%%%%%%%%%%%%%%%%%%%%%%%%%%%%%%%%%%%%%%%%%%%%%%%%%%%%%%%%%%%%%%%%%%%
% New Commands
%%%%%%%%%%%%%%%%%%%%%%%%%%%%%%%%%%%%%%%%%%%%%%%%%%%%%%%%%%%%%%%%%%%%%%%%%%%%%%%%%

% Pfeil in Zeichnungen
\newcommand{\pfeil} {Latex[scale width =0.8, scale length=1.2,bend]}

%Überschriften
\setcounter{secnumdepth}{2} % Zahlen nur bis section
\setcounter{tocdepth}{2}	% nur bis section in ToC
\renewcommand{\autodot}{}% Remove all end-of-counter dots
\AtBeginDocument{
  \renewcommand{\chapterautorefname}{Chapter}
  \renewcommand{\sectionautorefname}{Section}
  \renewcommand{\subsectionautorefname}{Subsection}
}

%Neue Textarten

\newcommand{\Name}[1]{\textsc{#1}}
\newcommand{\eigenname}[1]{"#1"}
\newcommand{\benennung}[1]{\textit{#1}}
\newcommand{\code}[1]{\texttt{#1}}
\newcommand{\zeichen}[1]{\(#1\)}

%linebreak nach \paragraph{}

\newcommand{\myparagraph}[1]{\paragraph{#1}\mbox{}\\}

%%%%%%%%%%%%%%%%%%%%%%%%%%%%%%%%%%%%%%%%%%%%%%%%%%%%%%%%%%%%%%%%%%%%%%%%%%%%%%%%%
%%%%%%%%%%%%%%%%%%%%%%%%%%%%%%%%%%%%%%%%%%%%%%%%%%%%%%%%%%%%%%%%%%%%%%%%%%%%%%%%%
%%%%%%%%%%%%%%%%%%%%%%%%%%%%%%%%%%%%%%%%%%%%%%%%%%%%%%%%%%%%%%%%%%%%%%%%%%%%%%%%%
\hyphenation{in-stan-ta-ne-ous}
\hyphenation{com-po-nents}
\hyphenation{bet-ween}


\begin{document}

\typeout{>>> TEXTWIDTH = \the\textwidth}
%=============================================================================%


\pagestyle{empty}



\includepdf[pages=-]{Deckblatt.pdf}



\begin{titlepage}
	\begin{center}


\begin{figure}[htbp]
	\centering
		
\includegraphics[width=75mm]{bilder/LSTM_Logo.pdf} 		
\end{figure}
		
		\vspace{25mm}
		 		
		\LARGE {\textbf{Direct numerical simulation \\ of turbulent flow \\ in rough pipes}}\\ [25mm]
		\huge {\textbf{Master's Thesis}} \\ [14mm]
		 by \\ [2mm]
		\LARGE \textbf{Jonas Tiefenthäler}\\[5mm]
		\LARGE Stud.reg.no. 23302719\\[5mm]
		\LARGE Erlangen, 31.07.2026\\
		\vfill
		\rule{15cm}{0.02mm}\\
		\Large Institute of Fluid Mechanics\\
		Faculty of Engineering\\
		Friedrich--Alexander--University Erlangen--N\"urnberg\\
		Prof.\ Dr.\ Philipp Schlatter\\

  		\end{center}
\end{titlepage}


  	\thispagestyle{empty}	
	\cleardoublepage

	
	\begin{tabular}[H]{ll}
	  Author:& Jonas Tiefenthäler \\
	  Student Registration Number: & 23302719 \\
	  Institute: & Institute of Fluid Mechanics\\
	  Head of Institute: & Prof.\ Dr.\ Techn. Philipp Schlatter\\ 
	  Group and Supervisor:& Prof.\ Dr.\ Techn. Philipp Schlatter \\
	  		
	  Assigned: & 01.02.2026 \\
	  Submitted: & 31.07.2026 \\
	  \\
	  LSTM Technical Report: & 18/2026
	\end{tabular}	



\thispagestyle{empty}
\cleardoublepage

\pagenumbering{Roman}
\setcounter{page}{1}
\chapter*{Declaration}

I hereby declare formally that I have developed and written the enclosed thesis entirely by myself and without help of third parties and have not used sources or means without declaration in the text. Any thoughts or quotations, inferred (literally or analogously) from the sources, are marked as such. This thesis was not submitted in the same or substantially similar version to any other authority to achieve an academic grade.

The Friedrich--Alexander--Universität (FAU), represented by the Institute of Fluid Mechanics (LSTM), is granted a simple, free of charge, temporally and geographically unrestricted right to use the results of the work, including any protective rights and rights of use, for the purposes of research and teaching.
\\ [35mm]
Erlangen, the 31.07.2026 \hspace{4cm} \rule{5cm}{0.02mm}

\chapter*{Abstract}
The behaviour of turbulent flow in rough pipes is studied using the \acrfull{dns}.
The Neko simulation tool is used, which is a high-order spectral element solver.
The rough surfaces are randomly generated using an algorithm, which takes a random distribution and power spectrum as input parameters.
The generated random surface is mapped onto the mesh, to create a body-fitted simulation mesh.

A total of nine simulations are performed.
Three different surface topographies are studied.
The three surfaces exhibit two different roughness heights.
Two of these surfaces have a reduced roughness height, with varying effective slopes.
All three surfaces are studied using three different Reynolds numbers.
The chosen Reynolds numbers are $5300, 11700$ and $19000$, corresponding to the canonical smooth pipe cases of the friction Reynolds numbers $180, 360$ and $550$.

Various data is presented.
This includes first- and second-order statistics, such as mean velocity profiles, total stress budgets and Reynolds stresses.
Furthermore, energy spectra are computed and presented.
Additionally, the development of viscous drag with respect to the roughness height and effective slope of the rough surface is examined.

The influence of the roughness height and effective slope of the surface on turbulent behaviour can be observed.
Furthermore, a clear and consistent impact on the energy spectra of the surfaces emerge.
A correlation is detected between the viscous drag, roughness height and effective slope.

\chapter*{Acknowledgment}

Firstly, I would like to thank Prof. Philipp Schlatter for supervising this thesis and for giving me the opportunity to write it.
I especially want to thank him, for spending his valuable time explaining me numerous aspects of turbulence statistics, the neko code and DNS in general.
This was invaluable to the success of this thesis.
I would also like to thank Dr. Siavash Toosi, for the countless questions he answered me during the last six months.
Especially our numerous discussions about the correct way to post-process rough pipes, helped significantly for the outcome of this work.
Furthermore, I would like to thank all the other students and employees at the LSTM for always making work enjoyable.


%=============================================================================%

%\thispagestyle{empty}



%=============================================================================%
\clearpage
\thispagestyle{empty}
\null
\clearpage
\thispagestyle{empty}
\null
\clearpage

\phantomsection
%\thispagestyle{empty}
\clearpage
\pagestyle{plain}




%=============================================================================%

\pagestyle{headings}


\addcontentsline{toc}{chapter}{Contents}


\tableofcontents

\listoffigures

\listoftables

%\printnomenclature
%=============================================================================%
%\addcontentsline{toc}{chapter}{Nomenklatur}
\newpage


%\thispagestyle{empty}

%\printnomenclature
\renewcommand\nomgroup[1]{%
  \item[\bfseries
  \ifstrequal{#1}{F}{Fluid properties}{%
  \ifstrequal{#1}{O}{Other symbols}{}}%
]}
%%%%%%%%%%%%%%%%%%%%%%%%%%%%%%%%%%%%%%%%%%%%%%%%%%%%%%%%%%%%%%
% Wenn Symbolverzeichnis automatisch erzeugt werdenn soll

\ifnomenklatur
%\ifenglish
%\printnomenclature\addcontentsline{toc}{chapter}{Notation}
%\else
%\printnomenclature\addcontentsline{toc}{chapter}{Nomenklatur}
%\fi

%%%%%%%%%%%%%%%%%%%%%%%%%%%%%%%%%%%%%%%%%%%%%%%%%%%%%%%%%%%%%%%%
\else
%% Wenn man das Symbolverzeichnis selber machen will
%\ifenglish
%\chapter*{Notation}\addcontentsline{toc}{chapter}{Notation}
%\else
\chapter*{List of Symbols}\addcontentsline{toc}{chapter}{List of Symbols}
%\fi
%
%
%
%
%%\thispagestyle{empty}
%\ifenglish
%\section*{Abbrevations}
%\else
%\section*{Abkürzungen}
%\fi
%
\nomenclature[F]{$\mu$}{Dynamic viscosity of the fluid}
\nomenclature[F]{$\rho$}{Density of the fluid}
\nomenclature[O]{$Re_b$}{Bulk Reynolds Number}
\nomenclature[O]{$Re_\tau$}{Friction Reynolds number}
\nomenclature[O]{C}{Mass flux}
\nomenclature[O]{$\alpha$}{Scaling factor to ensure constant mass flux}


\printnomenclature

% \begin{longtable}[t]{p{0.1\textwidth}p{0.1\textwidth}p{0.8\textwidth}}
% 	\textbf{Variable}  & \textbf{Einheit} & \textbf{Beschreibung}\\
% 	\midrule
% 	$c_0$ & $m/s$ & Schallgeschwindigkeit in Luft \\
% 	$\nabla$ & $1/m$ & Nabla-Operator \\
% 	$\mu$ & & Dynamic viscosity of the fluid\\
% 	$\rho$ & & Density of the fluid\\
% 	$\alpha$ & & Scaling factor to ensure constant mass flux\\
% 	$Re_\tau$ & & Friction Reynolds number\\
% 	$Re_b$ & & Bulk Reynolds number\\
% 	$C$ & & Mass flux\\

	
% 	%
% \end{longtable}


%\chapter*{Abkürzungsverzeichnis}\addcontentsline{toc}{chapter}{Abkürzungsverzeichnis}
%{\let\clearpage\relax \chapter*{List of Abbreviations}}\addcontentsline{toc}{chapter}{List of Abbreviations}

\newacronym{sem}{SEM}{Spectral Element}
\newacronym{cfd}{CFD}{computational fluid dynamics}
\newacronym{dns}{DNS}{Direct Numerical Simulation}
\newacronym{les}{LES}{Large Eddy Simulation}
\newacronym{nse}{NSE}{Navier-Stokes-Equations}
\newacronym{rans}{RANS}{Reynolds-averaged Navier-Stokes}
\newacronym{da}{DA}{Double-Averaged}
\newacronym{bc}{B.C}{Boundary Condition}
\newacronym{bcs}{B.Cs}{Boundary Conditions}
\newacronym{ic}{I.C}{Initial Condition}
\newacronym{ics}{I.Cs}{Initial Conditions}
\newacronym{gll}{GLL}{Gauss-Lobatto-Legrende}
\newacronym{fem}{FEM}{Finite-Element-Method}
\printglossary[type=\acronymtype]


% \begin{longtable}[t]{p{0.25\textwidth}p{0.8\textwidth}}
% 	\textbf{Abkürzung}  & \textbf{Bedeutung}\\
% 	\midrule
% 	NSE & Navier-Stokes Equations \\
	

	
	
% 	%
% \end{longtable}



\fi	
\cleardoublepage

\clearpage
\thispagestyle{empty}
\null
\clearpage
\thispagestyle{empty}
\null
%=============================================================================%

\pagenumbering{arabic}

\include{kapitel/einleitung}

%TODOS:

THEORY:
\begin{itemize}
    \item skin friction coefficient beschreiben
    \item noch was zu den roughness regimes schreiben
\end{itemize}
VALIDATION:
\begin{itemize}
    \item vlt. noch was schreiben, dass nur der vergleich der daten wichtig ist, nicht welche kurve was ist
\end{itemize}
RESULTS: 
\begin{itemize}
    \item checken ob erhöhte $v'v'$ werte an verdrängung durch rougness features liegt
    \item Schreiben dass längere dominante wavelenghts für geringere roughness heights daran liegt, dass die turbulenz weniger aufgebrochen wird.
    \item vlt noch was schreiben warum die energy peaks ca. bei den uu peaks sind
    \item noch schreiben, dass surfaces die turbulenz aufbrechen daher energy peaks an längeren/kürzeren lambda
    \item interpretation für höhere U+ bei y+ = 0 geben bei reynolds comp
\end{itemize}

\chapter{Introduction}\label{chap: Introduction}

Turbulent flows are ubiquitous in life.
Especially turbulent flows over a wall, so called wall-bounded turbulent flows, are encountered frequently in all kinds of technical applications ZITAT.
Around $50 \%$ of the energy usage in flows through pipes or the movement of cars or planes through air can be accounted to the energy dissipation of wall-bounded turbulent flow \citep{jimenezCascadesWallBoundedTurbulence2012}.
To increase the understanding of fluid dynamics, among others, the use of \acrfull{cfd} is a vital tool in fluid dynamics research.
Especially \acrfull{dns} and \acrfull{les} is a common choice among researchers, since it gives important insight inro the behaviour of turbulent flow ZITAT.
However, also the simulation \acrfull{rans} is a common choice, if the application and optimization of engineering tasks is the goal.

Turbulent flows over rough surfaces are mostly studied in the canonical channel case, e.g.\ by \citep{thakkarInvestigationTurbulentFlow2017,flackSkinFrictionMeasurements2020,kimTurbulenceStatisticsFully1987,jellyReynoldsDispersiveShear2018}, while the turbulent rough pipe has received less attention \citep{chanSystematicInvestigationRoughness2015,demaioDirectNumericalSimulation2023}.
Some studies used isolated roughness elements to find the impact of individual roughness features \citep{chatzikyriakouDNSTurbulentFlow2015}, while others tried to investigate the real world turbulence behaviour by using scanned surfaces of real objects \citep{thakkarSurfaceCorrelationsHydrodynamic2017,flackSkinFrictionMeasurements2020,frohnapfelFlowResistanceHeterogeneous2024}.\\

%\include{kapitel/theory_Michaela}

\chapter{Theory}\label{chap: Theory}


\section{Description of turbulent flows}\label{sec: Description of turbulent flows}

In turbulent flows, structures in a multitude of scales emerge.
The minimum scale of the elements is limited by the viscosity of the fluid, since the smallest scales are dissipated by viscosity and disappear.
Thus, the viscous units, or plus units are introduced. 

\subsection{Wall shear stress}\label{subsec: wall shear stress}

The wall shear stress, which is responsible for the drag, is defined as 
\begin{equation}
    \tau_w = \mu\frac{dU}{dy}
\end{equation}
and is solely based on viscosity (\cite{pope_turbulent_2015}).
The viscosity is dominant at the wall, hence the wall shear stress is the only stress active at the wall and it keeps being dominant in the near wall region.
With a growing distance to the wall the influence of viscosity decreases, while the importance of the Reynolds stresses, described in section \ref{subsec: Reynolds stress} increases.


\subsection{Friction velocity}\label{subsec: Friction velocity}

Using the wall shear stress and the density $\rho$ the friction velocity is defined as
\begin{equation}
    u_\tau = \sqrt{\frac{\tau_w}{\rho}}
\end{equation}
The friction velocity $u_\tau$ is a one of the viscous scales, which are used to a
\subsection{Reynolds stresses}\label{subsec: Reynolds stress}


%\section{Numerics}\label{sec: Numerics}

\section{Roughness}\label{sec: Roughness}

\subsection{Roughness parameters}\label{subsec: Roughness parameters}

To be able to describe rough surfaces, a multitude of roughness parameters exist, as described by \cite{chung_predicting_2021}.
The considered roughness parameters in this study are:

\paragraph{Average roughness height}

The average roughness height is the mean deviation  of the surface from the virtual origin $y=0$
\begin{equation}\label{eq: avg roughness height}
    S_a = \sum_{i,j}^{M,N} |h_{i,j}| / (MN)
\end{equation}

\paragraph{Root-mean-squeare roughness height}

\begin{equation}
    S_q = \sqrt{\sum_{i,j}^{M,N} h_{i,j}^{2}} / (MN)
\end{equation}

\paragraph{Skewness}

\begin{equation}
    Sk = \sum_{i,j}^{M,N} h^{3}_{i,j} /(MNS_{q}^{3})
\end{equation}

\paragraph{Maximum peak-to-valley height}
The maximum peak-to-valley height is the difference between the highest and lowest point of the rough surface
\begin{equation}
    S_{Z_{max}} = max(h_{i,j}) - min(h_{i,j})
\end{equation}

\paragraph{Effective slope}
The effective slope is the absolute gradient of the rough surface as described by \cite{napoli_effect_2008}.
\begin{equation}
    ES = \frac{1}{L_x L_y} \int_0^{L_x} \int_0^{L_y} \left| \frac{\mathrm{d}h}{\mathrm{d}x} \right|_{i,j} \mathrm{d}x \, \mathrm{d}y
\end{equation}
\paragraph{Mean peak-to-valley height}
The calculation of the mean peak to valley height is done according to \cite{thakkar_investigation_2017}.
The method is dividing the surface into tiles and calculating the difference between the highest and lowest point for each tile.
The mean of the differences results in  the mean peak-to-valley height $k$, which is also used as the measurement for roughness height in this thesis, according to \cite{de_maio_direct_2023}.
For this study the rough surface is divided into 25 rectangular tiles of equal size.
Equal size means in this context, that the same number of grid points are used for the length and width of the tiles.


\subsection{Virtual origin}\label{subsec: Virtual origin}

The virtual origin is a reference plane, which is taken as the basis for measuring the wall distance $y$ (\cite{chung_predicting_2021}).
In a smooth wall case, there is no virtual origin, since all data can be calculated from the wall itself, hence the point where the no-slip condition is active (\cite{chan_systematic_2015}), which leads to consistent wall distances for every point in the field.\\
Yet, in a rough wall case, this is not as straightforward.
Due to wall roughness, the distance from e.g. the pipe center to the wall varies continuously over the whole flow field.
To be able to compute consistent data, a virtual origin has to be defined. 
There are many different approaches to defining the virtual origin.
Notable methods are the usage of the mean momentum absorption plane by \cite{thom_momentum_1971} or the determination of the position of the mean streamwise velocity being zero, since it is assumed, that this is the equivalent position of the no slip condition in a smooth wall case.
A third approach is to take the mean radius of the rough surface for the virtual origin.\\
In the investigation by \cite{chan_systematic_2015}, all three of the approaches are compared. 
The approach of using the mean radius for the virtual origin seems to offer the best correlation between the considered smooth and rough cases in this study.
For that reason, this is the used method in this study.



%\include{kapitel/Methods_korrektur}
\chapter{Methods}\label{chap: Methods}

For this Thesis the high-order \acrfull{sem} code NEKO is used.
Neko is an open-source \acrfull{cfd} solver, that is based on the older Nek5000 \acrshort{cfd} solver done by \cite{fischer_nek500_nodate} in the nineties.
The portable framework Neko is written in modern Fortran in an object-oriented manner and is mainly focused on incompresible flow simulations \cite{jansson_neko_2024}.
Only \acrfull{dns} will be performed in this work, however \acrfull{les} models are also implemented in NEKO.

\section{Simulation setup}\label{sec: Simulation setup}

\section{Generating rough surfaces}\label{sec:generating rough surfaces}

\section{Postprocessing}\label{sec: Postprocessing}

In order to be able to work with the obtained data by the \acrshort{dns}, the data has to be postprocessed.
In this work, this is mainly done by using PySEMTools, which is a open source Python package, to postprocess \acrshort{sem} data on hexahedra high-order elements \cite{perez_pysemtools_2026}.\\
The typical workflow is as follows:
\begin{itemize}
    \item In case multiple files exist that contain statistics data, average them to obtain a single averaged field
    \item Interpolate the averaged data from the \acrshort{sem} mesh to user defined interpolation points
    \item Calculate the desired values based on the interpolated data
    \item Further average the data, such that 2D or 1D plots can be generated
\end{itemize}
To be able to process data of rough cases properly, various changes had to be done compared to a smooth case.

\subsection{Postprocessing rough cases}\label{subsec: Postprocessing rough cases}

The Postprocessing of rough \acrshort{dns} data is not as straightforward as for smooth cases.
Two of the several conceivable workflows will be described in this section.

\subsubsection{Mapping Interpolation points to match the rough surface}\label{subsubsec: Mapping interpolation points}
In this approach, the interpolation points are mapped to match the rough surface of the \acrshort{dns}.
Due to the matching field, the interpolation points can be interpolated onto the \acrshort{dns} field with all points being found by the interpolator.
The main issue with this approach is the fact, that the interpolation points, which contain the data that is used in the further process, are not uniformly distributed in the radial direction.
This will lead to the situation, that data points that are interpolated at different radii will be treated as if they were at a comparable position.

\subsubsection{Inverse mapping \acrshort{dns} data and interpolation points}

The second approach is to introduce a computational domain for the data interpolation.
To map the \acrshort{dns} data into the computational domain it is mapped to a smooth pipe.
This is achieved by applying the inverse of the surface function to the rough surface.
The inverse of the rough surface function can be expressed as
\begin{equation}
     surf_{inverse} = \frac{1.0}{(1.0 + \frac{surf_{change}}{pipe_R - r_{crit}})} * (mesh_{change} + \frac{r_{crit}}{pipe_R - r_{crit}} * surf_{change})
\end{equation}
One must account for the fact, that for performing correct computations for first- and second-order statistics, data at points with a uniform distance from the pipe center are necessary.




\chapter{Implementation}\label{chap: application}

The theoretical concepts and methods described in \autoref{chap: Theory} and \autoref{chap: Methods} must be incorporated into a \acrshort{dns} simulation setup. 
Furthermore, the \acrshort{dns} data must be postprocessed, to extract the relevant quantities presented in \autoref{chap: Results}.
In the following, first, the simulation setup used in this theses is described in \autoref{sec: Simulation setup}, followed by a brief description of the postprocessing procedure in \autoref{sec: post-processing}.
Finally, in \autoref{sec: performed simulations} the performed simulations and chosen parameters are described.

\section{Simulation setup}\label{sec: Simulation setup}
For this study, the canonical setup of the pipe is used.
In the pipe, azimuthal periodicity is given by default meaning that a smooth pipe setup is fully defined using the pipe length and the Reynolds-number \citep{yaoDirectNumericalSimulations2023}.
Additionally, a cylindrical coordinate system is employed, with the three spatial directions $r, \theta$ and $z$.
The terminology wall-normal is used for the $r$-coordinate, azimuthal is the $\theta$ direction and $z$ is the streamwise direction.
The cartesian coordinates used by Neko can be transformed into the cylindrical coordinate system by the relations $r = \sqrt{x^2 + y^2}$ and $\theta = \arctan\!\left(\frac{y}{x}\right)$.

The dimensions of the rough domain are $L_r \times L_\theta \times L_z = R \times 2\pi R \times 8\pi R$, with the radius $R$ is set to $1 \hat{=}$ unity. 
The smooth validation domains are shortened to $L_z = 4\pi R$ to save computation time.
$R$ denotes the mean radius of the pipe in rough setups, and is therefore the distance from the virtual origin to the pipe center.

For the \glsentryfull{bcs}, the Dirichlet no-slip \acrshort{bc} at the pipe wall is used.
This states that the fluid at the wall  moves with the same velocity as the wall itself.
In the current setup all velocities are set to zero at the wall: $u_i = 0 \ \text{for}\ y = 0$.
This also implies that the wall is impermeable, since the wall normal velocity $v$ is zero.
Since the streamwise direction is statistically homogeneous in fully developed turbulent pipe flow \citep{moinDIRECTNUMERICALSIMULATION1998}, a periodic boundary condition is used for the inflow and outflow of the pipe.
This means that the outflow velocities are set equal to the inflow velocities: $u_{i_{l = 0}} = u_{i_{l=L_z}}$.
This implies that fluid leaving the domain at the outflow reenters the domain at the inflow.

It is important to use appropriate \glsentryfull{ics} to support a fast transition to turbulence.
This is favourable since it saves computational resources.
This is achieved by using a laminar velocity profile as a basis, which is overlaid by large- and small-scale disturbances.
The laminar velocity profile is defined as
\begin{equation}\label{eq: laminar velocity profile}
    u_r = 0 \quad u_\theta = 0 \quad u_z = 2(1-r^2)
\end{equation}
\begin{figure}
    \centering
    \includegraphics[width = 0.5\textwidth]{plots/Methods/Initial_condition-cropped.pdf}
    \caption[Visualisation of initial condition]{Visualisation of initial condition. Blue colours denote lower velocities, red denotes higher velocities.}
    \label{fig: initial condition}
\end{figure}
with $r$ denoting the radius ranging from $0$ to $1$.
The used profile represents the bulk velocity $U_b = 1$ in streamwise direction.
The perturbations are defined as a combination of multiple $\sin$ and $\cos$ waves of various magnitudes.
Large- and small-scale disturbances are intended to disrupt the laminar profile and support a fast transition to turbulence.
The initial velocity field is taken from the Neko example of the turbulent channel and converted to polar coordinates, and altered, such that the divergence-free criterion $\nabla u_i = 0$ is met, which is necessary for the conservation of mass \citep{popeTurbulentFlows2015}.
A visualisation of the initial velocity magnitude is depicted in \autoref{fig: initial condition}.
It can be seen, that the \acrshort{ic} is not physical, due to a presence of singularity at the pipe centre.
This is however unavoidable, since it was chosen to use the principle of sinusoidal disturbances in order to keep it in line with the \acrshort{ic} in the Neko channel example, but with a met divergence free criterion in polar coordinates.

The bulk flow rate $U_b$ is defined in the streamwise $z$ direction as unity ($U_b=1$) and is kept constant throughout the simulation.
To achieve the desired bulk velocity, the flow is driven by a pressure gradient $\Pi$.
To maintain a constant flow rate, the pressure gradient $\Pi$ is varying over time to adjust for changes in the flow conditions, e.g., when turbulence develops.
Using \autoref{eq: viscous scaling units} it is also possible to compute $Re_\tau$ from the pressure gradient $\Pi$  as described by~\cite{demaioDirectNumericalSimulation2023}
\begin{equation} \label{eq: REtau from forcing}
    Re_\tau = \sqrt{\frac{\Pi}{2}}Re_b \quad \text{since} \quad \tau_w = \frac{R}{2}\Pi 
\end{equation}

As mentioned in \autoref{chap: Introduction}, past numerical studies of turbulent flows over rough surfaces frequently used the \acrshort{ibm} method \citep{demaioDirectNumericalSimulation2023,busseReynoldsnumberDependenceNearwall2017,thakkarSurfaceCorrelationsHydrodynamic2017,frohnapfelFlowResistanceHeterogeneous2024}.
For the \acrshort{ibm}, a uniform mesh is constructed, which is slightly larger than the domain of interest.
Within the uniform mesh, the nodes that lie on the boundary between the desired shape and the flow domain are identified.
At each of the boundary nodes an additional forcing term is added, which enforces the \glsentryfull{bcs} \citep{verziccoImmersedBoundaryMethods2023}.
When comparing \acrshort{ibm} with body-fitted mesh topologies, both advantages and disadvantages emerge.
An advantage of the \acrshort{ibm} method is the significantly easier mesh topoglogy, which simplifies mesh generation.
Furthermore, because of the orthogonal mesh alignment, the computation effort per node is significantly reduced for the \acrshort{ibm} compared to a body-fitted mesh.
In contrast, the required number of grid points is higher for the \acrshort{ibm} method, with a difference with increasing Reynolds number.
The ratio of necessary grid-points for the \acrshort{ibm} to a body-fitted mesh can be estimated as $N_{IBM}/N_{BF}\approx Re^{0.9}$ \citep{verziccoImmersedBoundaryMethods2023}.
This results in a lower computation cost for the \acrshort{ibm} method for lower Reynolds numbers, due to the lower computation effort per node, but a higher computation cost for higher Reynolds numbers.
In this thesis a body fitted-mesh is chosen in order to provide data for wall roughness resolved using body-fitted meshes.

The simulation mesh must contain enough grid points to be able to correctly resolve the physics of turbulence.
Since the turbulent scales decrease near the wall and the viscous sublayer must to be resolved, it is necessary to refine the grid resolution near the wall.
Vice versa, the mesh can be coarsened towards the pipe center, without compromising the accuracy of the results.
\begin{figure}
    \centering
    \includegraphics[width = 0.5\textwidth]{Methods/mesh_19000Reb_crossSection.pdf}
    \caption[Visualisation of simulation mesh]{Exemplary visualisation of simulation mesh. Structured mesh topology in the outer part, unstructured mesh topology in the inner part.}
    \label{fig: simulation mesh}
\end{figure}
The simulation mesh is generated using Gmsh \citep{Gmsh}.
\autoref{fig: simulation mesh} shows the frontal view on one plane of an exemplary simulation mesh.
All the other meshes used in this work have a very similar topography.
It can be seen that the mesh consists of two sections.
The outer part is generated using a structured mesh topology, while the inner part uses an unstructured mesh.
This is done to avoid heavily distorted elements and singularity at the mesh center.
The size of the structured region is defined by an inner radius, which establishes the boundary between the structured and unstructured sections.
Increased resolution in the near-wall section is achieved through a geometric progression of the element size in the structured region of the mesh.
The number of elements in the azimuthal direction, $N_\theta$, also implicitly defines the number of elements in the unstructured part, since the size or number of mesh elements cannot be explicitly defined in the unstructured mesh topography.
All simulation meshes consist purely of quad elements.
To determine reasonable mesh resolutions, a mesh convergence study is performed in \autoref{subsec: grid refinement study}.
The resulting mesh parameters from this study are listed in \autoref{tab: parameters chosen meshes}.

For the time discretization, described in \autoref{subsec: time discretization}, the \acrshort{cfl}-number of 0.5 is chosen throughout this study, eventhough safe \acrshort{cfl} number range until $\approx 0.7$.
This choice is made, since the \acrshort{cfl}-numer varies throughout the simulation in a range of $\pm 0.1$.

\subsection{Applying rough surfaces onto the mesh}\label{subsec: Generating rough mesh}

The algorithm described in \autoref{subsec: Surface generation algorithm} yields the height values for a plane surface.
This surface has to be mapped onto the pipe in such a way, that the same generated surface can be mapped onto different meshes with varying amounts of grid points.
%This leads to the fact, that the straightforward method, i.e. applying the height values of the generated surface points directly onto the pipe mesh is not applied in this work.
The used approach relies on the fact that every surface topography can be decomposed into a multitude of $\sin$ and $\cos$ waves using a Fourier transform \citep{perez-rafolsGeneratingRandomlyRough2019}.
The decomposed values are the Fourier coefficients, which can be seen as a weight given for each mode.
By recombining the extracted modes, the surface can be recreated independently of the simulation mesh.
The Fourier coefficients are extracted by taking the real and imaginary part  of the Fourier transform of the generated rough surface, such that $A(m,n) = \operatorname{Re}\left(\hat{Z}_h\right)$ and $B(m,n) = \operatorname{Im}\left(\hat{Z}_h\right)$.

This decomposition and recreation is only exact, if an infinite number of modes is used.
%This is however only relevant for regular surfaces.
%Since the surfaces generated by the algorithm described in section \ref{subsec: Surface generation algorithm} are defined with a finite amount of modes (that is chosen by the user), this amount of modes is sufficient for the recreation of the surface.
With these parameters, the height function for the rough surface at each individual point can be defined as
\begin{equation}\label{eq: height function}
    k(\theta,z) = \sum_{m,n}^{M,N} A(m,n)\cos(phase)- B(m,n)\sin(phase)
\end{equation}
with
\begin{equation}\label{eq: phase equation}
    phase = 2\pi\left(\frac{uz}{L_z} + \frac{v\theta}{L_{\theta}}\right)
\end{equation}
with $A(m,n)$ and $B(m,n)$ representing the Fourier coefficients, $m$ and $n$ the Fourier mode and $z,\ \theta$ the streamwise and azimuthal coordinate of the respective point.
$M$ and $N$ denote the total number of Fourier modes that are taken into account.
The Fourier transform is performed using the Numpy python library, the result is ordered such that the first half of the array $A[1:M/2]$ contains the positive frequency terms, while the latter half $A[M/2+1:]$ contains the negative frequency terms, which follows the~\cite{NumpyFT}.
To get wave-numbers which are signed correctly, i.e.\ positive, $m$ and $n$ are adjusted accordingly, which gives the parameters $u$ and $v$.

The total number of Fourier modes $M$ and $N$ changes the shape of the surface significantly.
A higher number of Fourier modes also includes higher wave-numbers, which leads to surfaces with smaller roughness features. 
This however, resuts in a higher grid resolution requirement, since it shows that at least $12$ grid points per roughness feature are necessary, to sufficiently resolve the shape of the roughness feature \citep{busseDirectNumericalSimulation2015}.

The height function in \autoref{eq: height function} yields the recreation of the rough surface, which is applied on grid points at the wall.
However, a body fitted mesh is desired, which leads to the necessity, to create a scaling function for the grid points away from the wall.
The roughness function should be applied onto the mesh in the near wall region of the domain, while the inner region stays unaffected \citep{chanSystematicInvestigationRoughness2015}.
The used function for this purpose is
\begin{equation}\label{eq: roughness scaling function}
    r_{shift} = r + \frac{r - r_{crit}}{R - r_{crit}} * k(\theta,z) 
\end{equation}
with $r_{shift}$ representing the amount of necessary radial shift of the respective point, $r_{crit}$ the minimal radius, at which the grid points are shifted, $R$ the maximum radius of the smooth pipe mesh and $r$ the radial position of the individual point to be shifted.
This function scales the magnitude of shift applied to each grid point linearly between the outer pipe radius of the original smooth pipe and $r_{crit}$, from one to zero.
This method results in a body fitted mesh, where the boundaries of the mesh represent the roughness features.
Furthermore, the near wall mesh refinement is kept.

\begin{figure}
    \centering
    \includegraphics[width = 0.5\textwidth]{Methods/deformed_Mesh-cropped.pdf}
    \caption[Body fitted rough mesh]{Mesh deformed due to roughness}
    \label{fig: deformed mesh}
\end{figure}
The points are only shifted in radial direction.
Since the GLL points are distributed over the mesh before the roughness is applied, the GLL distribution is not exact after the rough deformation.
This effect is assumed to be negligable, since the shortening or stretching of the elements is of small magnitude.

\section{Postprocessing}\label{sec: post-processing}
\fxbmark{Diese section noch nicht}
In order to be able to work with the obtained data by the \acrshort{dns}, a post-processing of the data is necessary.
The \acrshort{dns} data is distributed on the \acrshort{gll} points and non uniformly distributed in space.
The post-processing is done, to map the data onto user defined points, since this simplifies all further processing steps.
All data is presented in a wall normal direction and after post-processing, the \acrshort{dns} data is mapped onto wall normal planes.

In this work, this is done using PySEMTools, which is an open source Python package, to post-process \acrshort{sem} data on hexahedral high-order elements \citep{perezPySEMToolsLibraryPostprocessing2026}.
Using PySEMTools it is possible to obtain fluid statistics from multiple instantaneous fields, which will be done in \autoref{chap: Results} for the energy spectra.
Additionally, Neko has the capability to produce statistics files, in which the data is averaged in time and, if chosen, in one or more spatial directions of choice during the runtime of the simulation.
These statistics files are used to obtain the relevant data for the first- and second order statistics plots.
In the present thesis, no spatial averaging is performed by Neko, so the statistics files contain 3D data that is only averaged in time.

The typical steps for post-processing are as follows:
\begin{itemize}
    \item In case multiple files exist that contain statistics data, average them to obtain a single time averaged field
    \item Interpolate the averaged data from the \acrshort{sem} mesh to user defined interpolation points
    \item Calculate the desired values, e.g.\ velocity profiles and turbulent budgets, based on the interpolated data
    \item Further average the 3D data in space, if desired
\end{itemize}
Special measures have to be taken to process rough data correctly.
The post-processing of rough data is described in the following section. 
The goal is to give a brief overview and explain the workflow in principle, thus some details are neglected.
\subsection{Postprocessing rough cases}\label{subsec: Postprocessing rough cases}
The post-processing of rough \acrshort{dns} data is not as straightforward as for smooth cases.
Since the virtual origin is not at the wall at most positions, the interpolation point-cloud  must account for that.


The approach used in the present thesis is to map the \acrshort{dns} domain, as well as the interpolation points into a computational domain.
In the computational domain, the simulation domain is mapped from the rough shape into a smooth pipe, while the interpolation points are mapped such that they are at the respective positions of concentric circles in the physical domain.
For the rough simulation domain, this is achieved by applying the inverse of the \autoref{eq: roughness scaling function} onto the rough mesh using the height functions \autoref{eq: height function} and \autoref{eq: phase equation}.
The inverse of \autoref{eq: roughness scaling function} is
\begin{equation}\label{eq: inverse height function}
     {r_{shift}}^{-1} = \frac{1.0}{\left(R + \frac{k(\theta, z)}{R - r_{crit}}\right)} * \left(r + \frac{r_{crit}}{R - r_{crit}} * k(\theta, z)\right)
\end{equation}

\begin{figure}
    \centering
    \begin{subfigure}{0.45\textwidth}
        \subcaption{}
        \includegraphics[width = \textwidth]{Methods/averaged_field.png}
        \label{subfig: averaged field}
    \end{subfigure}
    \hfill
    \begin{subfigure}{0.45\textwidth}
        \subcaption{}
        \includegraphics[width = \textwidth]{Methods/smooth_Averaged_field.png}
        \label{subfig: smooth averaged field}
    \end{subfigure}
    \plotverticalspacing
    \caption[Visualisation of the smoothing algorithm used in the postprocessing]{Visualization of the result of the smoothing algorithm. Color depicting the velocity magnitude of a time averaged field.~(a) Rough pipe domain, (b) Rough domain mapped to smooth pipe}
    \label{fig: comparison of smooth and normal averaged field}
\end{figure}
The results are shown from an exemplary case in \autoref{fig: comparison of smooth and normal averaged field}.
It can be observed, that the rough pipe is mapped back to the original radius of $R = 1$.
Yet, the data is scaled properly, such that the structure of the flow remains preserved.

To obtain correct statistics, the data should be interpolated on concentric circles in the physical domain.
\begin{figure}
    \centering
    \begin{subfigure}{0.325\textwidth}
        \subcaption{}
        \input{plots/Methods/interpolation_point_placement_1.pgf}
        \label{subfig: PP points at beginning}
    \end{subfigure}
    \hfill
    \begin{subfigure}{0.325\textwidth}
        \subcaption{}
        \input{plots/Methods/interpolation_point_placement_2.pgf}
        \label{subfig: PP points after mapping to smooth}
    \end{subfigure}
    \hfill
    \begin{subfigure}{0.325\textwidth}
        \subcaption{}
        \input{plots/Methods/interpolation_point_placement_3.pgf}
        \label{subfig: PP points after inverse maping}
    \end{subfigure}
    \plotverticalspacing
    \caption[Placement of data points and interpolation points while mapping in postprocessing]{Visualization of points mapping in post-processing. Dots \dnsdot~\ represent \acrshort{dns} data points, circles  \interpcircle~\ interpolation points, the  black line \wallline~\ the pipe wall and the grey area \domainsquare~\ the simulation domiain (a) Default point placement. (b) DNS domain mapped into the smooth computational domain, as depicted in \autoref{subfig: smooth averaged field}. (c) Interpolation points mapped by the inverse of the rough surface, such that they lie on the same respective concentric positions in the computational domain}
    \label{fig: Visualization of the PP points}
\end{figure} 
If the smooth averaged field would be interpolated onto concentric interpolation points, the generated data would be in a body fitted manner.
To achieve data extraction in a concentric way, the interpolation points have to be deformed as well. 
By applying \autoref{eq: inverse height function} onto the interpolation points as done for the rough pipe, the interpolation points are deformed in such a way, that they lie in the computational domain in the same positions, as they would be in the physical domain for a concentric data extraction. This principle is shown in \autoref{fig: Visualization of the PP points}.

\section{Performed simulations}\label{sec: performed simulations}

To improve the understanding of the behaviour of turbulence over rough surfaces, a multitude of simulations are performed in the present thesis.
The rough surfaces as well as the Reynolds numbers are altered, to gain a comprehensive insight into the roughness effects in turbulent pipe flows.
A total number of nine simulations are performed.
Three different surfaces are considered, and all three surfaces are studied using three different Reynolds numbers.
The chosen bulk Reynolds numbers $Re_b$ are $5300,\ 11700$ and $19000$.
Those were chosen, since they are the Reynolds numbers for the typical friction Reynolds numbers of past studies in smooth pipes of $Re_\tau = 180,\ 360$ and $550$ \citep{yaoDirectNumericalSimulations2023,elkhouryDirectNumericalSimulation2013,kimProgressPipeChannel2012}. 
The considered surfaces are all based on the study by~\cite{demaioDirectNumericalSimulation2023}.
It is chosen to use the typical bulk Reynolds-numbers rather than match the friction Reynolds-numbers, since $Re_\tau$ does not describe the turbulence behaviour as it does in smooth cases.
This is due to the fact, that it is not only dependent on the viscous drag, but also on the pressure drag in rough conditions.
This makes it also dependent on the surface shape and not only on the separation of scales, hence the viscosity, like it is in smooth cases.

\begin{table}
    \centering
    \caption[Surface parameters of simulation surfaces]{Surface parameters of simulation surfaces. $S1$ showing the greatest roughness height, $S2$ reduced roughness height with similar effective slope, $S3$ is $S1$ scaled down such that the roughness height of $S2$ is matched}
    \begin{tabular}{c c c c c c c}
        \toprule
                & $S_a$     & $S_q$     & $Sk$      & $S_{z_{max}}$ & $ES$      & $S_{z_{5 \times 5}}$  \\
        \midrule
        $S1$    & $0.031$   & $0.040$   & $-0.479$  & $0.306$       & $0.261$   & $0.228$               \\
        $S2$    & $0.021$   & $0.026$   & $-0.500$  & $0.230$       & $0.256$   & $0.171$               \\
        $S3$    & $0.024$   & $0.030$   & $-0.479$  & $0.229$       & $0.196$   & $0.171$               \\
        \bottomrule
    \end{tabular}
    \label{tab: surface parameters of production runs}
\end{table}
In \autoref{tab: surface parameters of production runs} the considered surface parameters for the three surfaces are listed.
The surface referred to as S1 resembles the grit-blasted surface of~\cite{demaioDirectNumericalSimulation2023} as a basis.
The $S1$ surface is not exactly the same, but an approximation.
The approximation is done, by creating a surface, that shows the same surface parameters, as the surface studied by~\cite{demaioDirectNumericalSimulation2023}.
It is created using the algorithm described in \autoref{subsec: Surface generation algorithm}.

Additionally, the surfaces $S2$ and $S3$ show a decreased surface height $k$ compared to $S1$.
$S2$ keeps the effective slope, despite having a decreased roughness height $k$.
\begin{figure}
    \centering
    \begin{subfigure}{0.8\textwidth}
        \centering
        \subcaption{}
        \subimport{plots/Methods/}{surface_height_distribution_piro.pgf}
        \label{subfig: surface height map piro}
    \end{subfigure}
    \\[-20pt]
    \begin{subfigure}{0.8\textwidth}
        \centering
        \subcaption{}
        \subimport{plots/Methods/}{surface_height_distribution_kDown.pgf}
        \label{subfig: surface height map kDown}
    \end{subfigure}
    \\[-20pt]
    \begin{subfigure}{0.8\textwidth}
        \centering
        \subcaption{}
        \subimport{plots/Methods/}{surface_height_distribution_scaled.pgf}
        \label{subfig: surface height map scaled}
    \end{subfigure}
    \begin{subfigure}{0.8\textwidth}
        \centering
        \hspace{0.5cm}
        \subimport{plots/Methods/}{heightmap_colorbar.pgf}
    \end{subfigure}
    \caption[Height maps of simulation surfaces]{Height maps of simulation surfaces. Red donting surface retractions, green denoting surface elevations.~(a) $S1$, (b) $S2$, (c) $S3$}
    \label{fig: surface height maps}
\end{figure}

$S3$ is a scaled-down version of $S1$, which results also in a decreased effective slope.
\autoref{fig: surface height maps} shows the height maps for all three surfaces.
Red and green shadings depict extreme values, while yellow shadings denote the mean height.
It can be seen, that $S1$ in \autoref{subfig: surface height map piro} shows the largest deviations from the mean height.
Furthermore, $S2$ in \autoref{subfig: surface height map kDown} exhibits more, but less pronounced roughness features compared to $S1$. 
$S3$ in \autoref{subfig: surface height map scaled} shows the same height map as $S1$ but with overall damped values.

It is chosen to use the same physical surfaces for all Reynolds numbers.
This leads to changing non-dimensionalised roughness heights $k^+$ between the Reynolds numbers.
This is done, since a fixed $h^+$ would necessitate the generation of new surface for every surface-Reynolds number combination.
However, due to the randomness and irregularity of the surfaces, this would lead to different physical shapes of the roughness features, even though the surface parameters remain the same. 
The similarity of flow conditions for surfaces with the same statistical surface parameters is not given, as discussed by~\cite{forooghiUniversalRoughnessCorrelation2017,chungPredictingDragRough2021}.
To mitigate this issue, the same surfaces were chosen for all Reynolds numbers.
This leads to different roughness heights in viscous length scales $k^+$ between the Reynolds numbers.

The cases with lowered roughness height $k$ show both the same roughness height.
It is chosen to be $\approx 75\%$ of the roughness height of $S1$.
This choice is made, since it is the lowest surface height with matched a $Es$, that can be generated using an appropriate number of Fourier modes that still enables the resolution requirement by~\cite{busseDirectNumericalSimulation2015}.
It can be seen, that all the surface parameters, that are comparable between the surfaces are within the margin of $0.05$, which is considered as sufficient.
The only exception is skewness, for which it proved difficult to find a good fit with the used algorithm.



\chapter{Validation} \label{chap: Validation}

\section{Smooth wall pipe}\label{sec: Smooth wall pipe}


\section{Sinusoidal wall pipe}\label{sec: Sinusoidal wall pipe}

\subsection{Postprocessing rough Simulations}\label{subsec: postprocess rough Simulations}

In order to obtain correct data from the rough simulations, the friction calculation in the postprocessing must be expanded.
The friction calculation is changed, such it is able to account for the rough surface.\\
DESCRIBE NEW POSTPROCESSING 
\\
To be able to validate the friction results of the \Name{Python}-postprocessing script, $\tau_w$ is computed by using the force scaling algorithm of \Name{Neko}.
The force scaling factor $alpha$, that is used in \Name{Neko} to ensure a constant mass flux, can also be used to compute $Re_\tau$ with
\begin{equation}
    Re_\tau = \sqrt{\alpha*0.5}*Re_b
\end{equation}
while $\alpha$ is computed using
\begin{equation}
    \alpha = \frac{C_{\text{target}} - C_{\text{current}}}{C_{\text{base}}}
\end{equation}
with $C_{\text{target}}$ denoting the desired mass flux, $C_{\text{current}}$ denoting the mass flux at the current time step and $C_{\text{base}}$ denoting the base mass flux, at the start of the simulation CHECKEN OB DAS mit C base STIMMT!!!.



\chapter{Results}\label{chap: Results}

\subsection{5700 Reb}
\begin{figure}
    \centering
    \begin{subfigure}[t]{0.495\textwidth}
        \centering
        \subcaption{}
        %\input{plots/Results/}
        \input{plots/Results/5700/5700Reb_comparison_uplus.pgf}
    \end{subfigure}
    \hfill
    \begin{subfigure}[t]{0.495\textwidth}
        \centering
        \subcaption{}
        \input{plots/Results/5700/5700Reb_comparison_RS.pgf}
    \end{subfigure}
    \hfill
    \begin{subfigure}[t]{0.495\textwidth}
        \centering
        \subcaption{}
        \input{plots/Results/5700/5700Reb_comparison_totalStress.pgf}
    \end{subfigure}
    \hfill
    \begin{subfigure}[t]{0.49\textwidth}
        \centering
        \subcaption{}
        \input{plots/Results/5700/5700Reb_velocityDefect.pgf}        
    \end{subfigure}
\end{figure}


\chapter{Conclusion}\label{chap: Conclusion}

FUTUREWORK
\begin{itemize}
    \item A test on the influence on the simulation results is beyond the scope of this work. 

\end{itemize}


%%%%%%%%%%%%%%%%%%%%%%%%%%%%%%%%%%%%%%%%%%%%%%%%%%%%%%%%%%%%%%%%%%%%%%%%%%%%%
%%%%%%%%%%%%%%%%%%%%%%%%%%%%%%%%%%%%%%%%%%%%%%%%%%%%%%%%%%%%%%%%%%%%%%%%%%%%%
%%%%%%%%%%%%%%%%%%%%%%%%%%%%%%%%%%%%%%%%%%%%%%%%%%%%%%%%%%%%%%%%%%%%%%%%%%%%%
%%%%%%%------Literaturverzeichnis----------%%%%%%%%%%%%%%%%%%%%%%%%%%%%%%%%%%
%%%%%%%%%%%%%%%%%%%%%%%%%%%%%%%%%%%%%%%%%%%%%%%%%%%%%%%%%%%%%%%%%%%%%%%%%%%%%
%%%%%%%%%%%%%%%%%%%%%%%%%%%%%%%%%%%%%%%%%%%%%%%%%%%%%%%%%%%%%%%%%%%%%%%%%%%%%
%%%%%%%%%%%%%%%%%%%%%%%%%%%%%%%%%%%%%%%%%%%%%%%%%%%%%%%%%%%%%%%%%%%%%%%%%%%%%

%\singlespacing
%\onehalfspacing

%\bibliographystyle{giubplain}
%\bibliographystyle{natdin}
%\bibliographystyle{plainnat_mod}
%\bibliographystyle{natdin}
%\bibliographystyle{dinat}


%%%%%%%%\bibliographystyle{unsrtdin_mod}
%\bibliographystyle{natdin}
\bibliographystyle{abbrvnat}
\bibliography{Bibliography}


%\bibliographystyle{natdin}
%\bibliographystyle{munich}



\newpage

\begin{appendix}                    % Anhang einfügen

\let\oldaddcontentsline\addcontentsline
\renewcommand{\addcontentsline}[3]{%
  \ifstrequal{#2}{chapter}{\oldaddcontentsline{#1}{#2}{#3}}{}%
}

%\chapter{Appendix}\label{chap: Appendix}

\chapter{Derivation of dispersive stress}\label{apx: derivation of dispersive stress}

This appendix derives the dispersive stress tensor introduced in \autoref{eq: dispersive stress tensor}.
The notation follows \autoref{par: dispersive stresses}.
An overline $\overline{ \phantom{u} }$ denotes a time average, angle brackets $\langle \phantom{u} \rangle$ denote a spatial average over the wall-parallel plane, a prime $ \phantom{u}'$ a fluctuation in time about the time average, and a tilde $\tilde{ \ }$ a fluctuation in space about the spatial average, i.e.\ $\tilde{u}_i = \overline{u_i} - \langle \overline{u_i}\rangle$.

\section{Analytical derivation}

The Reynolds stresses are defined as the time average of the product of the temporal velocity fluctuations
\begin{equation}\label{eq: reynolds stresses}
    R_{ij} = \overline{(u_i - \overline{u_i})(u_j - \overline{u_j})}
    = \overline{u_i' u_j'}
\end{equation}
The dispersive stress tensor follows from applying the same operation to the time-averaged velocity $\overline{u_i}$, but with the average taken in space rather than in time:
\begin{align}
    D_{ij}  &= \langle (\overline{u_i} - \langle \overline{u_i}\rangle)(\overline{u_j} - \langle \overline{u_j}\rangle) \rangle \notag\\
            &= \langle \overline{u_i}\,\overline{u_j} - \overline{u_i}\langle\overline{u_j}\rangle - \langle \overline{u_i}\rangle\overline{u_j} + \langle\overline{u_i}\rangle\langle\overline{u_j}\rangle \rangle \notag\\
    %\intertext{Since $\langle\overline{u_i}\rangle$ and $\langle\overline{u_j}\rangle$ no longer vary in space, they can be pulled out of the outer spatial average term by term:}
            &= \langle \overline{u_i}\,\overline{u_j} \rangle - \langle\overline{u_i}\rangle\langle\overline{u_j}\rangle - \langle\overline{u_i}\rangle\langle\overline{u_j}\rangle + \langle\overline{u_i}\rangle\langle\overline{u_j}\rangle \notag\\
            &= \langle \overline{u_i}\,\overline{u_j} \rangle - \langle\overline{u_i}\rangle\langle\overline{u_j}\rangle \notag\\
    %\intertext{Substituting $\overline{u_i} = \langle\overline{u_i}\rangle + \tilde{u}_i$ (and analogously for $\overline{u_j}$) and using that the spatial average of a spatial fluctuation vanishes, $\langle\tilde{u}_i\rangle = 0$, leaves only the cross term of the fluctuations:}
            &= \langle \tilde{u}_i\,\tilde{u}_j \rangle \label{eq: dispersive stress}
\end{align}
which recovers the definition of the dispersive stress tensor given in \autoref{eq: dispersive stress tensor}.

% \section{Evaluation in the post-processing script}

% The post-processing script computes the pointwise Reynolds stress $R_{ij} = \overline{u_i'u_j'}$ at every spatial location directly from the time-resolved velocity field, before any spatial averaging is applied.
% To recover the correctly double-averaged Reynolds stress $\langle R_{ij}\rangle$, the term $\langle \overline{u_i}\,\overline{u_j}\rangle - \langle\overline{u_i}\rangle\langle\overline{u_j}\rangle$ is added and subtracted, which makes the dispersive stress from \autoref{eq: dispersive stress} appear explicitly:
% \begin{align}
%     \langle R_{ij}\rangle   &= \langle R_{ij} + \overline{u_i}\,\overline{u_j} \rangle - \langle \overline{u_i} \rangle \langle\overline{u_j}\rangle \notag\\
%                             &= \langle \overline{u_i' u_j'} \rangle + \langle \overline{u_i}\,\overline{u_j} \rangle - \langle \overline{u_i} \rangle \langle\overline{u_j}\rangle \notag
%     \intertext{using \autoref{eq: dispersive stress} for the last two terms:}
%                             &= \langle \overline{u_i' u_j'} \rangle + \langle \tilde{u}_i\,\tilde{u}_j\rangle \notag\\
%                             &= \langle \overline{u_i' u_j'} \rangle + D_{ij} \label{eq: reynolds stress in PP}
% \end{align}
% i.e.\ the double-averaged Reynolds stress splits into the spatially averaged, pointwise Reynolds stress and the dispersive stress tensor.
% This is the identity used by the post-processing script: $\overline{u_i'u_j'}$ and $\overline{u_i}$ are computed directly from the simulation output, and $\langle\overline{u_i'u_j'}\rangle$ and $D_{ij}$ are then obtained from the spatial average and \autoref{eq: dispersive stress}, respectively.

% --- Original version, kept for reference ---
% The Reynolds stresses are defined as:
% \begin{equation}\label{eq: reynolds stresses}
%     R_{ij} = \overline{(u_i - \overline{u_i})(u_j - \overline{u_j})}
%     = \overline{u_i' u_j'}
% \end{equation}
%
% DErivation of Dispersive stresses.
% The * superscripts denoting spatial fluctiuations.
% \begin{equation}\label{eq: dispersive stress}
%     \begin{aligned}
%         D_{ij}  &= \langle (\overline{u_i} - \langle \overline{u_i}\rangle)(\overline{u_j} - \langle \overline{u_j}\rangle) \rangle\\
%                 &= \langle \overline{u_i}\,\overline{u_j} - \overline{u_i}\langle\overline{u_j}\rangle - \langle \overline{u_i}\rangle\overline{u_j} + \langle\overline{u_i}\rangle\langle\overline{u_j}\rangle \rangle\\
%                 &= \langle \overline{u_i}\,\overline{u_j} \rangle - \langle \overline{u_i}\langle\overline{u_j}\rangle \rangle - \langle \langle \overline{u_i}\rangle\overline{u_j} \rangle + \langle \langle\overline{u_i}\rangle\langle\overline{u_j}\rangle\rangle \\
%                 &= \langle \overline{u_i}\rangle \langle \overline{u_j}\rangle + \langle \overline{u_i}^* \, \overline{u_j}^*\rangle - \langle\overline{u_i}\rangle\langle\overline{u_j}\rangle - \langle\overline{u_i}\rangle\langle\overline{u_j}\rangle + \langle\overline{u_i}\rangle\langle\overline{u_j}\rangle \\
%                 &= \langle \overline{u_i}^* \, \overline{u_j}^*\rangle
%     \end{aligned}
% \end{equation}
%
% Derivation of the calculations of Reynolds and dispersive stresses, as done in the postprocessing script.
% \begin{equation}\label{eq: reynolds stress in PP}
%     \begin{aligned}
%         \langle R_{ij}\rangle   &= \langle R_{ij} + \overline{u_i}\,\overline{u_j} \rangle - \langle \overline{u_i} \rangle \langle\overline{u_j}\rangle  \\
%                                 &= \langle \overline{u_i' u_j'} + \overline{u_i}\,\overline{u_j} \rangle - \langle \overline{u_i} \rangle \langle\overline{u_j}\rangle \\
%                                 &= \langle \overline{u_i' u_j'} \rangle + \langle \overline{u_i}\,\overline{u_j} \rangle - \langle \overline{u_i} \rangle \langle\overline{u_j}\rangle \\
%                                 &= \langle \overline{u_i' u_j'} \rangle  + \langle \overline{u_i}\rangle \langle\overline{u_j}\rangle + \langle \overline{u_i}^* \,\overline{u_j}^*\rangle - \langle \overline{u_i} \rangle \langle\overline{u_j}\rangle \\
%                                 &= \langle \overline{u_i' u_j'} \rangle + \langle \overline{u_i}^* \,\overline{u_j}^*\rangle
%     \end{aligned}
% \end{equation}
%
% \begin{equation} \label{eq: dispersive stress in PP}
%     \begin{aligned}
%         D_{ij}  &= \langle\overline{u_i}\,\overline{u_j}\rangle - \langle\overline{u_i}\rangle\langle\overline{u_j}\rangle\\
%                 &= \langle \overline{u_i}\rangle \langle \overline{u_j}\rangle + \langle \overline{u_i}^* \, \overline{u_j}^*\rangle - \langle\overline{u_i}\rangle\langle\overline{u_j}\rangle\\
%                 &= \langle \overline{u_i}^* \, \overline{u_j}^*\rangle
%     \end{aligned}
% \end{equation}

\chapter{Energy spectra}\label{apx: Energy budgets}
\section{Energy spectra for $Re_b = 11700$}
\begin{figure}[h!]
    \centering
    \begin{subfigure}{0.495\textwidth}
        \subcaption{}
        \subimport{plots/Results/11700/}{11700_spectra_piro.pgf}
        \label{subfig: 11700 spectra piro}
    \end{subfigure}
    \hfill
        \begin{subfigure}{0.495\textwidth}
            \subcaption{}
            \subimport{plots/Results/11700/}{11700_spectra_kDown.pgf}
            \label{subfig: 11700 spectra kDown}
    \end{subfigure}
    \\[-20pt]
    \begin{subfigure}{0.495\textwidth}
        \subcaption{}
        \subimport{plots/Results/11700/}{11700_spectra_scaled.pgf}
        \label{subfig: 11700 spectra scaled}
    \end{subfigure}
    \centering
    \begin{subfigure}{\textwidth}
        \centering
        \hspace{0.9cm}
        \subimport{plots/Results/11700/}{3DSpectra_colorbar.pgf}
    \end{subfigure}
    \caption[Streamwise pre-multiplied energy spectra at $Re_b = 11700$]{Streamwise pre-multiplied Energy spectra at $Re_b = 11700$. Wall-normal position plotted over wavelength. Colour denoting energy content, with blue indicating low energy content, and yellow high energy content. Red marker indicates the position of maximal energy. Contour lines start at 0.25 with intervals of 0.5.~(a) $S1$.~(b) $S2$.~(c) $S3$}
    \label{fig: energy spectra 11700}
\end{figure}

\section{Energy spectra for $Re_b = 19000$}
\begin{figure}[h!]
    \centering
    \begin{subfigure}{0.495\textwidth}
        \subcaption{}
        \subimport{plots/Results/19000/}{19000_spectra_piro.pgf}
        \label{subfig: 19000 spectra piro}
    \end{subfigure}
    \hfill
        \begin{subfigure}{0.495\textwidth}
            \subcaption{}
            \subimport{plots/Results/19000/}{19000_spectra_kDown.pgf}
            \label{subfig: 19000 spectra kDown}
    \end{subfigure}
    \\[-20pt]
    \begin{subfigure}{0.495\textwidth}
        \subcaption{}
        \subimport{plots/Results/19000/}{19000_spectra_scaled.pgf}
        \label{subfig: 19000 spectra scaled}
    \end{subfigure}
    \centering
    \begin{subfigure}{\textwidth}
        \centering
        \hspace{0.9cm}
        \subimport{plots/Results/19000/}{3DSpectra_colorbar.pgf}
    \end{subfigure}
    \caption[Streamwise pre-multiplied energy spectra at $Re_b = 19000$]{Streamwise pre-multiplied Energy spectra at $Re_b = 19000$. Wall-normal position plotted over wavelength. Colour denoting energy content, with blue indicating low energy content, and yellow high energy content. Red marker indicates the position of maximal energy. Contour lines start at 0.25 with intervals of 0.5.~(a) $S1$.~(b) $S2$.~(c) $S3$}
    \label{fig: energy spectra 19000}
\end{figure}
\newpage
\section{Energy spectra for $S2$}

\begin{figure}[h!]
    \centering
    \begin{subfigure}{0.495\textwidth}
        \subcaption{}
        \subimport{plots/Results/comp_Reynolds/kDown/}{5300_spectra_kDown.pgf}
        \label{subfig: 5300 spectra kDown reyndolsComp}
    \end{subfigure}
    \hfill
        \begin{subfigure}{0.495\textwidth}
            \subcaption{}
            \subimport{plots/Results/comp_Reynolds/kDown/}{11700_spectra_kDown.pgf}
            \label{subfig: 11700 spectra kDown reynoldsComp}
    \end{subfigure}
    \\[-20pt]
    \begin{subfigure}{0.495\textwidth}
        \subcaption{}
        \subimport{plots/Results/comp_Reynolds/kDown/}{19000_spectra_kDown.pgf}
        \label{subfig: 19000 spectra kDwon reynoldsComp}
    \end{subfigure}
    \centering
    \begin{subfigure}{\textwidth}
        \centering
        \hspace{0.9cm}
        \subimport{plots/Results/comp_Reynolds/kDown/}{3DSpectra_colorbar.pgf}
    \end{subfigure}
    \caption[Streamwise pre-multiplied energy spectra for $S2$]{Streamwise pre-multiplied Energy spectra for $S2$ at the three Reynolds numbers. Wall-normal position plotted over wavelength. Colour denoting energy content, with blue indicating low energy content, and yellow high energy content. Red marker indicates the position of maximal energy. Contour lines start at 0.25 with intervals of 0.5.~(a) $Re_b = 5300$.~(b) $Re_b = 11700$.~(c) $Re_b = 19000$}
    \label{fig: energy spectra S2}
\end{figure}
\newpage
\section{Energy spectra for $S3$}
\begin{figure}[h!]
    \centering
    \begin{subfigure}{0.495\textwidth}
        \subcaption{}
        \subimport{plots/Results/comp_Reynolds/scaled/}{5300_spectra_scaled.pgf}
        \label{subfig: 5300 spectra scaled reyndolsComp}
    \end{subfigure}
    \hfill
        \begin{subfigure}{0.495\textwidth}
            \subcaption{}
            \subimport{plots/Results/comp_Reynolds/scaled/}{11700_spectra_scaled.pgf}
            \label{subfig: 11700 spectra scaled reynoldsComp}
    \end{subfigure}
    \\[-20pt]
    \begin{subfigure}{0.495\textwidth}
        \subcaption{}
        \subimport{plots/Results/comp_Reynolds/scaled/}{19000_spectra_scaled.pgf}
        \label{subfig: 19000 spectra scaled reynoldsComp}
    \end{subfigure}
    \centering
    \begin{subfigure}{\textwidth}
        \centering
        \hspace{0.9cm}
        \subimport{plots/Results/comp_Reynolds/scaled/}{3DSpectra_colorbar.pgf}
    \end{subfigure}
    \caption[Streamwise pre-multiplied energy spectra for $S3$]{Streamwise pre-multiplied Energy spectra for $S3$ at the three Reynolds numbers. Wall-normal position plotted over wavelength. Colour denoting energy content, with blue indicating low energy content, and yellow high energy content. Red marker indicates the position of maximal energy. Contour lines start at 0.25 with intervals of 0.5.~(a) $Re_b = 5300$.~(b) $Re_b = 11700$.~(c) $Re_b = 19000$}
    \label{fig: energy spectra S3}
\end{figure}



\chapter{Development of viscous drag with roughness height and streamwise slope}\label{apx: viscous drag to surface height}

\section{Surfaces $S2$ and $S3$ at $Re_b = 5300$}
\begin{figure}[h!]
    \centering
    \subimport{plots/Results/5300/}{5300Reb_kDowntauWall_over_height.pgf}
    \caption{Roughness height and viscous drag over streamwise position for  $S2$ at $5300 Re_b$}
\end{figure}
\begin{figure}[h!]
    \centering
    \hspace{0.9cm}
    \subimport{plots/Results/5300/}{5300Reb_kDown_tauWall_over_height_bySlope.pgf}
    \caption[Scatterplot of viscous drag over physical roughness height at $Re_b = 5300$ for $S2$]{Scatter plot of nondimensionalised viscous drag over roughness height. Colour denoting local streamwise slope. Data shown for $S2$ at $5300 Re_b$}
\end{figure}
\begin{figure}[h!]
    \centering
    \subimport{plots/Results/5300/}{5300Reb_scaledtauWall_over_height.pgf}
    \caption{Roughness height and viscous drag over streamwise position for  $S3$ at $5300 Re_b$}
\end{figure}
\begin{figure}[h!]
    \centering
    \hspace{0.9cm}
    \subimport{plots/Results/5300/}{5300Reb_scaled_tauWall_over_height_bySlope.pgf}
    \caption[Scatterplot of viscous drag over physical roughness height at $Re_b = 5300$ for $S3$]{Scatter plot of nondimensionalised viscous drag over roughness height. Colour denoting local streamwise slope. Data shown for $S3$ at $5300 Re_b$}
\end{figure}
\clearpage
\section{Surfaces $S2$ and $S3$ at $Re_b = 11700$}
\begin{figure}[h!]
    \centering
    \subimport{plots/Results/11700/}{11700Reb_kDowntauWall_over_height.pgf}
    \caption{Roughness height and viscous drag over streamwise position for  $S2$ at $11700 Re_b$}
\end{figure}
\begin{figure}[h!]
    \centering
    \hspace{0.9cm}
    \subimport{plots/Results/11700/}{11700Reb_kDown_tauWall_over_height_bySlope.pgf}
    \caption[Scatterplot of viscous drag over physical roughness height at $Re_b = 11700$ for $S2$]{Scatter plot of nondimensionalised viscous drag over roughness height. Colour denoting local streamwise slope. Data shown for $S2$ at $1700 Re_b$}
\end{figure}
\begin{figure}[h!]
    \centering
    \subimport{plots/Results/11700/}{11700Reb_scaledtauWall_over_height.pgf}
    \caption{Roughness height and viscous drag over streamwise position for  $S3$ at $11700 Re_b$}
\end{figure}
\begin{figure}[h!]
    \centering
    \hspace{0.9cm}
    \subimport{plots/Results/11700/}{11700Reb_scaled_tauWall_over_height_bySlope.pgf}
    \caption[Scatterplot of viscous drag over physical roughness height at $Re_b = 11700$ for $S3$]{Scatter plot of nondimensionalised viscous drag over roughness height. Colour denoting local streamwise slope. Data shown for $S3$ at $11700 Re_b$}
\end{figure}

\clearpage
\section{Surfaces $S2$ and $S3$ at $Re_b = 19000$}

\begin{figure}[h!]
    \centering
    \subimport{plots/Results/19000/}{19000Reb_kDowntauWall_over_height.pgf}
    \caption{Roughness height and viscous drag over streamwise position for  $S2$ at $19000 Re_b$}
\end{figure}
\begin{figure}[h!]
    \centering
    \hspace{0.9cm}
    \subimport{plots/Results/19000/}{19000Reb_kDown_tauWall_over_height_bySlope.pgf}
    \caption[Scatterplot of viscous drag over physical roughness height at $Re_b = 19000$ for $S2$]{Scatter plot of nondimensionalised viscous drag over roughness height. Colour denoting local streamwise slope. Data shown for $S2$ at $19000 Re_b$}
\end{figure}
\begin{figure}[h!]
    \centering
    \subimport{plots/Results/19000/}{19000Reb_scaledtauWall_over_height.pgf}
    \caption{Roughness height and viscous drag over streamwise position for  $S3$ at $19000 Re_b$}
\end{figure}
\begin{figure}[h!]
    \centering
    \hspace{0.9cm}
    \subimport{plots/Results/19000/}{19000Reb_scaled_tauWall_over_height_bySlope.pgf}
    \caption[Scatterplot of viscous drag over physical roughness height at $Re_b = 19000$ for $S3$]{Scatter plot of nondimensionalised viscous drag over roughness height. Colour denoting local streamwise slope. Data shown for $S3$ at $19000 Re_b$}
\end{figure}

\let\addcontentsline\oldaddcontentsline
\end{appendix}


\end{document}

